% Cover
\newpage
\thispagestyle{empty}

% Background
\begin{tikzpicture}[remember picture,overlay]
    \fill[biru](current page.north west) rectangle ([yshift=-297mm]current page.north east);
    \fill[white] (current page.north west) rectangle ([yshift=-45mm]current page.north east);
\end{tikzpicture}

% Logo ITS
\begin{tikzpicture}[remember picture, overlay]
    \node[anchor=south west, inner sep=0] at (-1,0.1) {\includegraphics[width=4cm]{Figs/Logo ITS.png}};
\end{tikzpicture}

\setlength{\parindent}{0pt}
\vspace{10ex}
{\fontspec{Trebuchet MS Bold}{
  \fontsize{14pt}{14pt}\selectfont \textcolor{white}{TUGAS AKHIR - SF4801}}
  }

\vspace{12ex}

\doublespacing
\begin{large}
  {\fontspec{Trebuchet MS Bold}{
  \fontsize{18pt}{16pt}\selectfont \textcolor{white}{Pemetaan Sebaran dan Pemodelan Bawah Permukaan Potensi Panas Bumi Berdasarkan Data Landsat 8 dan Metode Gravitasi}}
  }
\end{large}

\vspace{15ex}

{\fontsize{14pt}{14pt}\textcolor{white}{\fontspec{Trebuchet MS Bold}Surya Anoraga Justitia Yusman}\\
\fontsize{14pt}{14pt}\textcolor{white}{\fontspec{Trebuchet MS}NRP 01111940000019}
}

\vspace{10ex}

{\fontsize{14pt}{14pt}\textcolor{white}{\fontspec{Trebuchet MS}Dosen Pembimbing}\\
\fontsize{14pt}{14pt}\textcolor{white}{\fontspec{Trebuchet MS Bold}Dr.rer.nat. Eko Minarto}\\
\fontsize{14pt}{14pt}\textcolor{white}{\fontspec{Trebuchet MS}NIP 19750205 199903 1 004}
}

\vspace{15ex}

\onehalfspacing
{\fontsize{12pt}{14pt}\fontspec{Trebuchet MS Bold}\textcolor{white}{DEPARTEMEN FISIKA} \\
\fontsize{12pt}{14pt}\fontspec{Trebuchet MS}\textcolor{white}
{Fakultas Sains dan Analitika Data \\
Institut Teknologi Sepuluh Nopember \\
Surabaya \\
2023}
}